% Response letter to reviewers -- a clean portable starter.
\documentclass[11pt]{article}
\usepackage[utf8]{inputenc}
\usepackage[T1]{fontenc}
\usepackage[a4paper,margin=25mm]{geometry}
\usepackage{microtype}
\usepackage{xcolor}
\usepackage{tcolorbox}
\usepackage{enumitem}
\usepackage[hidelinks]{hyperref}
\definecolor{commentframe}{HTML}{64748B}\definecolor{commentfill}{HTML}{F1F5F9}
\definecolor{responseframe}{HTML}{0F766E}\definecolor{responsefill}{HTML}{F0FDFA}
\newtcolorbox{reviewercomment}{colback=commentfill,colframe=commentframe,boxrule=.6pt,arc=2pt,left=8pt,right=8pt,top=6pt,bottom=6pt}
\newtcolorbox{authorresponse}{colback=responsefill,colframe=responseframe,boxrule=.6pt,arc=2pt,left=8pt,right=8pt,top=6pt,bottom=6pt}
\newcommand{\reviewcomment}[1]{\begin{reviewercomment}\textbf{Reviewer comment. }#1\end{reviewercomment}}
\newcommand{\reviewresponse}[1]{\begin{authorresponse}\textbf{Response. }#1\end{authorresponse}}
\setlist{nosep}\title{Response to Reviewers' Comments}\author{Maya Chen and Alex Morgan}\date{29 August 2026}
\begin{document}
\maketitle
\noindent\textbf{Manuscript:} ``Reliable Data Pipelines for Small Research Teams''\\
\textbf{Journal:} Journal of Open Methods\\\textbf{Manuscript ID:} JOM-2026-014
\section*{Cover note to the editor}
Thank you for the careful and constructive review of our manuscript. We have
addressed every point below and marked the corresponding changes in the
revised manuscript. Page and section references refer to the revised version.
\section*{Reviewer 1}
\reviewcomment{The evaluation does not explain how the three data sources were selected, so it is difficult to assess whether the results generalise beyond the example workflow.}
\reviewresponse{We added a source-selection subsection in Section 3.1. It explains the inclusion criteria, reports the number of candidate sources screened, and adds a sensitivity analysis for the smallest source.}
\reviewcomment{The reproducibility claim would be stronger if the configuration values used in the experiments were visible in the paper.}
\reviewresponse{We added Table 2 with the complete configuration, including the validation window, retry policy, and random seed. The companion repository now contains the exact machine-readable configuration file.}
\section*{Reviewer 2}
\reviewcomment{Please distinguish failures caused by malformed input from failures caused by unavailable network services.}
\reviewresponse{We agree. Section 4.2 now reports the two failure classes separately, and Figure 3 uses different symbols for input validation and service errors.}
\section*{Summary of changes}
\begin{itemize}\item Added source-selection criteria and sensitivity analysis.\item Added a complete experimental configuration table.\item Split malformed-input and network-service failure reporting.\end{itemize}
\end{document}
